\chapter{Nonlinear residual observables}
\label{mht:ch:nonlinear-neumann}

The residual algebra has no higher products, but recovering a residual
field is a different operation. A matrix-valued connection must first be
put into Neumann gauge. Its normalized frame solves a linear equation
whose coefficient is the connection. The solution depends nonlinearly
on that coefficient. The first jet of this frame gives the residual
matrix and all its invariant polynomials. An explicit sphere family
will have zero linear projection and a nonzero quadratic invariant.

\section{The normalized frame}

Retain $C,o,j,E,G,f_0,\eta_0$ and $\kappa$ from
Chapter~\ref{mht:ch:polar-neumann}. Fix $N\geq1$ and let $t$ be a formal
parameter of degree zero. A formal curve connection is
\[
 a\in t\bigl(E^1\otimes\operatorname{Mat}_N(\mathbb C)\bigr)[[t]].
\]
All series in this section are $t$-adically complete. The operator $G$
acts entrywise on matrix-valued one-forms.

\begin{definition}\label{mht:def:normalized-frame}
Let $L_a(F)=aF$. The normalized frame and the residual matrix of $a$ are
\begin{align}
 H(a)&=(1+GL_a)^{-1}I_N
       =I_N-Ga+G(aGa)-G(aG(aGa))+\cdots,\label{mht:eq:frame}\\
 \mathcal J(a)&=-jH(a).\label{mht:eq:nonlinear-residual}
\end{align}
\end{definition}

Each coefficient uses finitely many applications of $G$ and
multiplications of a one-form by a function. The recursion and the
normalization of $G$ show that $H(a)$ is the unique formal solution of
\begin{equation}\label{mht:eq:frame-equation}
 \bar\partial H=-aH,\qquad H(o)=I_N.
\end{equation}
Its constant term is $I_N$, so it is formally invertible.
On a complex curve there is no further $(0,2)$ integrability equation.

A formal Neumann gauge is a matrix
$k\in\operatorname{GL}_N(C^\infty(C)[[t]])$ whose reduction modulo $t$
is a constant invertible matrix and whose first jet satisfies $j(k)=0$.
It acts by
\[
 a^k=k^{-1}ak+k^{-1}\bar\partial k.
\]
The gauge is based if $k(o)=I_N$. The condition on its constant term
ensures that $a^k$ again vanishes modulo $t$.

\begin{theorem}\label{mht:thm:neumann-normal-form}
For every formal curve connection, the matrix
\begin{equation}\label{mht:eq:normal-gauge}
 g(a)=H(a)\exp(f_0\mathcal J(a))
\end{equation}
is a based Neumann gauge and satisfies
$a^{g(a)}=\eta_0\mathcal J(a)$. Under a Neumann gauge,
\begin{equation}\label{mht:eq:frame-equivariance}
 H(a^k)=k^{-1}H(a)k(o),\qquad
 \mathcal J(a^k)=k(o)^{-1}\mathcal J(a)k(o).
\end{equation}
More precisely, gauges with $a^k=b$ correspond bijectively to
$K\in\operatorname{GL}_N(\mathbb C[[t]])$ satisfying
$\mathcal J(a)K=K\mathcal J(b)$, by
\begin{equation}\label{mht:eq:gauge-morphisms}
 k=H(a)KH(b)^{-1},\qquad K=k(o).
\end{equation}
\end{theorem}

\begin{proof}
The frame equation gives $a^{H(a)}=0$, and
$0^{\exp(f_0\mathcal J)}=\eta_0\mathcal J$.
The factors in \eqref{mht:eq:normal-gauge} have value $I_N$ at $o$ and
first jets $-\mathcal J$ and $\mathcal J$. Thus their product is a based
Neumann gauge.

The matrix $k^{-1}H(a)k(o)$ solves the normalized frame equation for
$a^k$. Uniqueness gives the first identity in
\eqref{mht:eq:frame-equivariance}. Taking its first jet, using
$j(k^{-1})=0$, gives the second.

If $a^k=b$, both $kH(b)$ and $H(a)$ are frames for $a$. Their ratio is
a global holomorphic matrix and hence is constant. This yields
\eqref{mht:eq:gauge-morphisms}. Conversely, that formula satisfies
$\bar\partial k=-ak+kb$, and
\[
 j(k)=-\mathcal J(a)K+K\mathcal J(b).
\]
It is Neumann precisely when the indicated intertwining relation holds.
\end{proof}

Every $J\in t\operatorname{Mat}_N(\mathbb C)[[t]]$ occurs, since
$H(\eta_0J)=\exp(-f_0J)$ and $\mathcal J(\eta_0J)=J$.
Thus the quotient by based gauges is exactly this formal matrix space.
Constant gauges then act by conjugation. The construction of $H$ and
$\mathcal J$ does not use $f_0$. That choice enters only the representative
$\eta_0\mathcal J$.

The theorem describes a neighborhood of the trivial holomorphic bundle.
It does not assert that every holomorphic bundle on the sphere is
trivial. In a small Sobolev neighborhood, the same construction is
convergent.

\begin{proposition}\label{mht:prop:analytic-frame}
Fix an integer $s\geq3$ and smooth metrics for the Sobolev norms on $C$.
For $a$ in a sufficiently small neighborhood of zero in
$H^s(\Omega_C^{0,1}\otimes\operatorname{Mat}_N(\mathbb C))$, the series
\eqref{mht:eq:frame} converges in $H^{s+1}$ to an invertible frame.
The maps $a\mapsto H(a)$ and $a\mapsto\mathcal J(a)$ are analytic.
Smooth inputs give smooth frames, also for smooth families parametrized
by compact subsets of a surface.
\end{proposition}

\begin{proof}
The Dolbeault solution operator gains one derivative. Subtracting its
value at $o$ preserves this bound, since evaluation is continuous on
$H^{s+1}$. Sobolev multiplication on the compact real surface gives
\[
 \|G(aF)\|_{H^{s+1}}
       \leq C_s\|a\|_{H^s}\|F\|_{H^{s+1}}.
\]
The Neumann series converges when $C_s\|a\|_{H^s}<1$.

After shrinking the neighborhood, Sobolev embedding makes $H-I_N$
uniformly small, so $H$ is invertible. First-jet evaluation is continuous
on $H^{s+1}$, proving analyticity of $\mathcal J$. The frame equation
bootstraps smooth inputs to smooth solutions. Differentiation of the
inverse operator with respect to surface parameters gives the same
estimates on each compact parameter set.
\end{proof}

Gauge equivariance continues to hold whenever both $a$ and $a^k$ belong
to this neighborhood. No smallness assumption on an unrelated gauge
orbit is implicit.

\section{The linear trace and the cyclic representative}

The first term of \eqref{mht:eq:nonlinear-residual} is
$p_1a=\kappa^{-1}\tau_E(a)$. The trace removes all its higher corrections,
although the matrix itself usually retains them.

\begin{proposition}\label{mht:prop:linear-first-trace}
Formally, and in the analytic neighborhood above,
\[
 \operatorname{tr}\mathcal J(a)=\kappa^{-1}\tau_E(\operatorname{tr}a).
\]
If the coefficient matrices of $a$ lie in one commutative subalgebra of
$\operatorname{Mat}_N(\mathbb C)$, then
$\mathcal J(a)=\kappa^{-1}\tau_E(a)$.
\end{proposition}

\begin{proof}
The formal or near-identity logarithm satisfies
\[
 \bar\partial\log\det H
 =\operatorname{tr}(H^{-1}\bar\partial H)=-\operatorname{tr}a.
\]
Its first jet is $-\operatorname{tr}\mathcal J(a)$.
The polar Stokes identity proves the first assertion.
In a commutative coefficient algebra, $\exp(-Ga)$ solves the normalized
frame equation. Taking its first jet proves the second.
\end{proof}

The cyclic homotopy of Theorem~\ref{mht:thm:cyclic-contraction} gives a
different gauge normalization. Seek $\xi\in t(E^0\otimes
\operatorname{Mat}_N(\mathbb C))[[t]]$ with $p_0(\xi)=0$, and expand
\[
 a^{e^\xi}=e^{-\operatorname{ad}\xi}a+
     \sum_{k\geq0}\frac{(-1)^k}{(k+1)!}
                   (\operatorname{ad}\xi)^k d_E\xi,
 \qquad \operatorname{ad}\xi(x)=[\xi,x].
\]
At order $n$, let $R_n$ denote the coefficient with the unknown
$\xi_n$ omitted. Set
\begin{equation}\label{mht:eq:cyclic-slice-recursion}
 \xi_n=-hR_n,\qquad J_n^{\mathrm{cyc}}=p_1R_n.
\end{equation}
The contraction identities give
$R_n+d_E\xi_n=\eta_0p_1R_n$ and $p_0(\xi_n)=0$, so this recursion
puts the connection in the form $\eta_0J^{\mathrm{cyc}}$.
It is unique: a first difference $\zeta$ of two normalized choices
satisfies $p_0\zeta=0$ and $hd_E\zeta=0$, hence $\zeta=0$.

For an input of field degree one, the first two homogeneous terms are
\begin{equation}\label{mht:eq:cyclic-first-terms}
 J_1^{\mathrm{cyc}}=p_1a,\qquad
 J_2^{\mathrm{cyc}}=-\frac12p_1[a,ha].
\end{equation}
Indeed, substituting $\xi_1=-ha$ in the quadratic coefficient leaves
$-\frac12[a,ha]-\frac12[\eta_0p_1a,ha]$ before projection.
The projection of the second commutator is zero, since
$p_1(\eta_0h\eta)=p_0(h\eta)=0$ entrywise.
Theorem~\ref{mht:thm:neumann-normal-form} makes
$J^{\mathrm{cyc}}$ conjugate to $\mathcal J(a)$.
Their invariant polynomials therefore agree.

\section{An exact noncommuting sphere family}

Take $N=2$, write $s_0=z\bar z$, and put
\[
 f=\frac{z^2}{(1+s_0)^2},\qquad
 g=\frac{\bar z}{1+s_0},\qquad
 A=E_{12},\quad B=E_{21},\quad
 M_0=[A,B]=\operatorname{diag}(1,-1).
\]
Here $E_{ij}$ are the elementary matrices. Both functions are globally
smooth: in the coordinate $w=1/z$ they become
$\bar w^2/(1+w\bar w)^2$ and $w/(1+w\bar w)$.
Their values and first holomorphic jets at $o$ vanish. Consequently
\[
 a=t(A\bar\partial f+B\bar\partial g)
\]
is exact in the linear Neumann complex, and $p_1a=0$.

\begin{theorem}\label{mht:thm:exact-sphere-residual}
For this family, as a formal identity and for sufficiently small
complex $t$,
\begin{equation}\label{mht:eq:exact-sphere-residual}
 \mathcal J(a)=-\frac{t^2}{3}M_0-\frac{13t^5}{315}B.
\end{equation}
It satisfies $\mathcal J(a)^2=(t^4/9)I_2$. For every $m\geq1$,
\begin{equation}\label{mht:eq:exact-sphere-invariants}
 \operatorname{tr}\mathcal J(a)^{2m+1}=0,\qquad
 \frac1{2m}\operatorname{tr}\mathcal J(a)^{2m}
       =\frac{t^{4m}}{m\,3^{2m}}.
\end{equation}
The constant Neumann gauge $K=I_2+(13/210)t^3B$ obeys
$K^{-1}\mathcal J(a)K=-(t^2/3)M_0$.
\end{theorem}

\begin{proof}
The normalized primitive
\[
 G(f\bar\partial g)=-\frac{z}{3(1+s_0)^3},
 \qquad f\bar\partial g=\frac{z^2}{(1+s_0)^4}\,d\bar z
\]
has first jet $-1/3$. These identities follow by differentiation,
smoothness in both charts, and the zero value at $o$.
The polar Stokes identity applied to $fg$ gives
$\tau_E(g\bar\partial f)=-\tau_E(f\bar\partial g)$.
The quadratic term
$-\kappa^{-1}\tau_E(aGa)$ of $\mathcal J(a)$ is therefore
$-t^2M_0/3$.

Write $H=\sum_{n\geq0}t^nH_n$, with $H_0=I_2$. The recursion
$\bar\partial H_n=-(a/t)H_{n-1}$ gives
\begin{align*}
 (H_2)_{11}&=\frac{z(3s_0+1)}{3(1+s_0)^3},\\
 (H_3)_{21}&=-\frac{s_0^4+4s_0^3+6s_0^2+2s_0}{6(1+s_0)^4},\\
 (H_4)_{11}&=\frac{z^2(15s_0^4+60s_0^3+90s_0^2+48s_0+8)}
                         {90(1+s_0)^6},\\
 (H_5)_{21}&=\frac{z(35s_0^4+140s_0^3+210s_0^2+126s_0+26)}
                         {630(1+s_0)^7}.
\end{align*}
Each right side is smooth in both charts and vanishes at $o$.
Differentiating, with
$\bar\partial f=-2z^3(1+s_0)^{-3}d\bar z$ and
$\bar\partial g=(1+s_0)^{-2}d\bar z$, verifies the recursion
and fixes each primitive uniquely. The last first jet is $13/315$.

To exclude all other orders, use rotations $z\mapsto e^{i\theta}z$.
The normalized solution operator commutes with their pullback, and
$f,g$ have weights $2,-1$. Since $A^2=B^2=0$, only alternating matrix
words occur. A word of length $2k$ has weight $k$, while words of length
$2k+1$ beginning with $A$ and $B$ have weights $k+2$ and $k-1$.

A smooth function can have nonzero first holomorphic jet at $o$ only
in weight one. The only surviving orders are therefore $n=2$ and the
$B$ word at $n=5$. This proves \eqref{mht:eq:exact-sphere-residual}
in every formal order.

Proposition~\ref{mht:prop:analytic-frame}
gives the analytic identity.

Now $M_0^2=I_2$, $M_0B+BM_0=0$, and $B^2=0$ give the square and all
trace identities. Finally $K^{-1}=I_2-(13/210)t^3B$ and
$[M_0,B]=-2B$ give the asserted conjugation.
\end{proof}

The first nonzero quadratic measurement is thus
$\frac12\operatorname{tr}\mathcal J(a)^2=t^4/9$.
The linear exactness of $a$ does not contradict this invariant:
a straight line of infinitesimal gauge variations need not be a finite
gauge orbit. A pure-gauge path has higher terms that enforce zero
residual.

The kinetic differential matters even for this small matrix example.
For comparison, take the Laurent complex
\[
 E_{\mathrm L}^0=\mathbb C[z,z^{-1}],\quad
 E_{\mathrm L}^1=E_{\mathrm L}^0\eta,\quad
 d_{\mathrm L}u=(z\partial_zu)\eta,\quad
 \tau_{\mathrm L}(u\eta)=[z^0]u.
\]
Here $\eta$ has degree one and square zero, and $[z^0]$ extracts the
constant coefficient. Its contraction projects to constant modes and
divides a nonzero mode by its integer exponent.

Set
\[
 a_{\mathrm L}=t\begin{pmatrix}0&z\\z^{-1}&0\end{pmatrix}\eta,\quad
 b(t)=\frac{\sqrt{1+4t^2}-1}{2},\quad r(t)=b(t)/t,\quad
 G_0=\begin{pmatrix}1&-rz\\rz^{-1}&1\end{pmatrix},
\]
using the formal square root with constant term one. Since
$b^2+b=t^2$,
\[
 z\partial_zG_0+(a_{\mathrm L}/\eta)G_0
       =G_0\operatorname{diag}(b,-b).
\]

Multiplying $G_0$ by $(1+r^2)^{-1/2}$ leaves the transformed connection
unchanged. Its logarithm is
$\arctan(r)\left(\begin{smallmatrix}0&-z\\z^{-1}&0\end{smallmatrix}\right)$,
which has zero constant Laurent coefficient. The resulting quadratic
invariant is
\[
 b(t)^2=t^4-2t^6+5t^8-14t^{10}+O(t^{12}).
\]
It belongs to this Laurent complex. The sphere invariant is
\eqref{mht:eq:exact-sphere-invariants}.

\section{Invariant functions and their smooth derivatives}

For each integer $r\geq1$, define
\begin{equation}\label{mht:eq:invariant-insertion}
 \mathcal O_r(a)=\frac1r\operatorname{tr}\mathcal J(a)^r.
\end{equation}
These functions are gauge invariant and satisfy
$\mathcal O_r(\eta_0J)=\operatorname{tr}(J^r)/r$.
Their Taylor coefficients are iterated Green expressions. For example,
superscripts denoting homogeneous field degree,
\begin{align*}
 \mathcal O_2^{[2]}&=\frac1{2\kappa^2}\operatorname{tr}\tau_E(a)^2,\\
 \mathcal O_2^{[3]}&=-\frac1{\kappa^2}
              \operatorname{tr}\bigl(\tau_E(a)\tau_E(aGa)\bigr),\\
 \mathcal O_2^{[4]}&=\frac1{\kappa^2}
              \operatorname{tr}\bigl(\tau_E(a)\tau_E(aG(aGa))\bigr)
              +\frac1{2\kappa^2}\operatorname{tr}\tau_E(aGa)^2.
\end{align*}
The polar trace acts entrywise. The final trace is a matrix trace.

\begin{proposition}\label{mht:prop:smooth-invariant-gradient}
The variation of $\mathcal O_r$ is represented by the smooth matrix
function
\begin{equation}\label{mht:eq:smooth-invariant-gradient}
 X_r(a)=\kappa^{-1}H(a)\mathcal J(a)^{r-1}H(a)^{-1},\qquad
 \delta\mathcal O_r=\tau_E\operatorname{tr}(X_r\,\delta a).
\end{equation}
It lies in $E^0\otimes\operatorname{Mat}_N(\mathbb C)$ and obeys
$\bar\partial X_r+[a,X_r]=0$.
\end{proposition}

\begin{proof}
Write $\delta H=HY$. Variation of the frame equation gives
\[
 \bar\partial Y=-H^{-1}\delta aH,\qquad Y(o)=0,\qquad
 Y=-G(H^{-1}\delta aH).
\]
Taking the first jet, using $H(o)=I_N$, yields
$\delta\mathcal J=\kappa^{-1}\tau_E(H^{-1}\delta aH)$.
Differentiate the invariant polynomial and use cyclicity of matrix
trace to obtain \eqref{mht:eq:smooth-invariant-gradient}.
The first jet of $X_r$ is
$\kappa^{-1}[-\mathcal J,\mathcal J^{r-1}]=0$.
Differentiating $H\mathcal J^{r-1}H^{-1}$ by $\bar\partial$ proves
covariant closedness.
\end{proof}

Let $U$ be a closed oriented surface with a smooth two-form $\rho$.
Let $a_C$ be a smoothly $U$-dependent curve connection in the formal or
analytic domain above. The functionals
\[
 I_r=\int_U\rho\,\mathcal O_r(a_C)
\]
have smooth representatives $\rho X_r$ in the conjugate antifield
component. They are local in $U$ and generally nonlocal in $C$.
The same construction applies with compact support on an open surface,
or with boundary conditions that annihilate the boundary variations.

For the Batalin--Vilkovisky (BV) extension, use the shifted space $L_U[1]$.
Its components of total form degrees $0,1,2,3$ are respectively the
ghost $c$, physical connection $\mathcal A$, connection antifield, and
ghost antifield. Their coordinate ghost numbers are $1,0,-1,-2$.
The pairing $\int_U\tau_E\operatorname{tr}$ pairs complementary total
degrees summing to three. After the shift, it has degree $-1$.
We use its degree-one antibracket on functionals whose differentials
admit smooth Hamiltonian representatives.

\begin{proposition}\label{mht:prop:classical-source-deformation}
Let $S_{\mathrm{BV}}$ be the minimal classical BV extension of
\eqref{mht:eq:open-action}, with field--antifield pairing induced by
$\int_U\tau_E\operatorname{tr}$.
Assume a boundary condition for which its classical master equation
holds. Then
\[
 \{S_{\mathrm{BV}},I_r\}=0,\qquad \{I_r,I_s\}=0.
\]
For any fixed finite set of degree-zero parameters $\sigma_r$, the action
$S_{\mathrm{BV}}+\sum_r\sigma_rI_r$ satisfies the classical master
equation. This remains true after completion in the ideal generated by
these sources.
\end{proposition}

\begin{proof}
On the physical curve connection, the minimal BV differential is its
infinitesimal gauge transformation with the ghost as parameter.
Gauge invariance of $\mathcal O_r$ gives the first identity, and
Proposition~\ref{mht:prop:smooth-invariant-gradient} supplies its
admissible smooth Hamiltonian representative. The $I_r$ depend only
on fields, so their mutual field--antifield contractions vanish.
Expansion of the master equation gives the assertion.
\end{proof}

One may also adjoin source antifields $\sigma_r^+$ of degree $-1$.
With antibracket convention $\{\sigma_r,\sigma_r^+\}=1$ and
$Q=\{S_{\mathrm{BV}}+\sum_r\sigma_rI_r,-\}$, the extended differential satisfies
$Q\sigma_r=0$ and $Q\sigma_r^+=I_r$.
This is a classical source deformation. It specifies no quantum
measure or coincident-point contraction.

\section{Direct sums and the closed-field comparison}

The frame equation respects block sums. Consequently the same
normalization of \eqref{mht:eq:invariant-insertion} at every rank gives
\[
 \mathcal O_r(a\oplus b)=\mathcal O_r(a)+\mathcal O_r(b).
\]
This property distinguishes single traces among homogeneous polynomial
families.

\begin{proposition}\label{mht:prop:additive-trace-family}
Let $F_N$ be a conjugation-invariant polynomial on
$\operatorname{Mat}_N(\mathbb C)$ for every $N\geq1$, homogeneous of one
fixed degree $d>0$. If
\[
 F_{N+M}(X\oplus Y)=F_N(X)+F_M(Y),
\]
then $F_N(X)=c_d\operatorname{tr}(X^d)$ with a scalar $c_d$ independent
of $N$. In degree zero, the additive constant families are $F_N=c_0N$.
\end{proposition}

\begin{proof}
Repeated direct-sum additivity gives
$F_N(\operatorname{diag}(x_1,\ldots,x_N))=\sum_iF_1(x_i)$.
Homogeneity gives $F_1(x)=c_dx^d$. Conjugation invariance then determines
$F_N$ on matrices with distinct eigenvalues, a Zariski dense set.
Equality holds as polynomials. In degree zero, additivity determines
the constant at rank $N$ as $N$ times the constant at rank one.
\end{proof}

For example, the direct-sum defect of
$u\operatorname{tr}(J^2)+v(\operatorname{tr}J)^2$ is
$2v\operatorname{tr}X\operatorname{tr}Y$.
Additivity forces $v=0$, while leaving $u$ undetermined.
It tests block-diagonal inputs and removes no off-diagonal fields at
the larger rank.

The invariant functions and their classical Hamiltonians are now
constructed. Identifying them as the image of geometric closed fields
requires another map, whose normalization and locality are independent
conditions.

\begin{problem}\label{mht:prob:closed-neumann-map}
Construct an action-derived map from an independently specified
geometric closed-field complex to the smooth Neumann insertions.
Determine its scalar normalization and whether its disk operation is
additive under block sums at fixed sources. Compare its full nonlinear
image with \eqref{mht:eq:invariant-insertion}, including the family of
Theorem~\ref{mht:thm:exact-sphere-residual}. The comparison must preserve
the chosen normal marking, boundary terms, two-source collisions, and
sewing. If it is required to be local in the curve variable, construct
that local comparison before using the global operator $G$.
\end{problem}
